\newpage

\section{Überschrift 1}
Im Inhaltsverzeichnis
\subsection{Überschrift 2}
Im Inhaltsverzeichnis
\subsubsection{Überschrift 3}
Im Inhaltsverzeichnis
\paragraph{Überschrift 4}
Nicht im Inhaltsverzeichnis
\subparagraph{Überschrift 5}
Nicht im Inhaltsverzeichnis

\subsection{Bilder}
Sind automatisch im Abbildungsverzeichnis
Nutzung von Bilder wie \autoref{fig:Hier die Beschriftung}
\flexFigure
{0.5}
{images/Bosch-Logo.pdf}
{Hier die Beschriftung}


\subsection{Tabellen}
Sind automatisch im Tabellenverzeichnis
Nutzung von Tabellen wie \autoref{tab:Farbaufteilung nach Lehrinhalt}
\flexTable
{|c|c|c|}
{
  \textbf{Farbe} & \textbf{Lehrinhalt} & \textbf{Klasse}
}
{
  Grün & Definitionen & alert alert-block alert-success \\
  Gelb & Beispiele & alert alert-block alert-warning \\
  Blau & Aufgaben & alert alert-block alert-info \\
}
{Farbaufteilung nach Lehrinhalt}

\subsection{Abkürzungen}
Erstes mal ausgeschrieben \ac{HTML} und danach abgekürzt \ac{HTML}.

\subsection{Quellcode}
Nutzung von Code wie \autoref{code:code/example.py}
\flexCode
{Python}
{Beschreibung}
{code/example.py}

\subsection{Zitieren}
% Direkte Zitate
„Software-Entwicklung umfasst alle Aspekte des Entwurfs, der Implementierung und des Tests von Software-Systemen“. \cite[23]{winniewski_grundlagenwissen_2024}
\\
% Indirekte Zitate
Die Software-Entwicklung beinhaltet verschiedene Phasen, von der Planung bis zur Wartung der Software \cite[45]{fuchs_disruption_2024}.
\\\textbf{Punktsetzung beachten}